\chapter{Introduction}

Managed AI coding agents are increasingly used as software-engineering assistants, but a thesis about such systems is weak if it only defines an evaluation framework. A framework can describe desired measurements, quality gates, and evidence paths, yet it does not demonstrate that the proposed infrastructure can execute on real repositories. The central problem of this work is therefore practical: the research project must turn local agent-evaluation ideas into reproducible engineering evidence without inventing measurements.

The objective is to build and execute a local evaluation harness around qyl and Arqio. qyl is treated as the research orchestration and OpenTelemetry surface, while Arqio is treated as the architecture-analysis component. The harness evaluates a small corpus of local and cloned repositories, records build and test outcomes, collects coverage when possible, checks semantic-convention compatibility, and produces thesis-ready tables from raw JSON artifacts. The expected measurable result is not agent correctness in the strong sense. Ground-truth labels are unset for the initial corpus, so the empirical claim is limited to harness feasibility, result reproducibility, coverage availability, telemetry availability, and label agreement for deterministic architecture classification.

This scope is intentionally narrower than a full vendor benchmark. A valid vendor comparison would require at least two configured managed-agent systems, preserved prompts, comparable tasks, and a sufficient sample size. If those conditions are not present, the thesis must state that managed-agent execution is unavailable. This is still an improvement over a purely conceptual paper because it replaces hypothetical evaluation with command logs, result JSON, generated tables, and failure reasons.

The engineering contribution is a repeatable local pipeline: target repository selection, restore/build/test execution, coverage collection, semantic-convention verification, Arqio-style classification, aggregation, Rego rubric input, and PDF generation. The thesis cites official documentation for OpenTelemetry, OPA/Rego, Vale, .NET testing, Coverlet, and GitHub Actions where those tools define the reproducible environment \cite{opentelemetry-dotnet,opentelemetry-semconv,opentelemetry-otlp,opa-rego,vale,dotnet-test,coverlet,github-actions-setup-dotnet}. qyl and Arqio are cited as local software artifacts because they are evaluated as private repositories rather than public scholarly sources \cite{qyl-repo,arqio-repo}.
